\subsubsection{坐标变换与流形上的 Sobolev 空间}\label{subsec:sobolev-coordinate-transformations-manifolds}

若 $\partial\Omega$ 光滑, 一个标准技巧是作局部坐标变换, 使 $\partial\Omega$ 成为坐标曲面. 这要求考察 Sobolev 空间在坐标变换下如何变化. 此外, 对 PDE 的边值问题, 需要考察定义在 $\partial\Omega$ 上的函数空间. 因而希望定义流形上的 Sobolev 空间. 当然, 这些空间将通过局部坐标卡来定义; 坐标变换下的行为同样是一个关键问题.

\begin{Definition}
\label{def:sobolev-k-diffeomorphism}
设 $\Omega,\Omega^{\prime}$ 为 $\mathbb{R}^m$ 中的开集. 若双射 $\bm{\Phi}:\Omega\to\Omega^{\prime}$ 满足以下条件, 则称其为 $k$ 阶微分同胚 $(k\geq1)$:\begin{enumerate}
\item $\bm{\Phi}$ 和 $\bm{\Phi}^{-1}$ 分别在 $\overline\Omega$ 和 $\overline{\Omega^{\prime}}$ 上连续,
\item 它们从 $1$ 阶直到 $k$ 阶的导数分别在 $\overline\Omega$ 和 $\overline{\Omega^{\prime}}$ 上有界且连续, 并且
\item 存在正常数 $c$ 和 $C$, 使在 $\Omega$ 中有 $c\leq|\det\nabla\bm{\Phi}|\leq C$.
\end{enumerate}
\end{Definition}

\begin{Definition}
\label{def:sobolev-diffeomorphism-pullbacks}
设 $\bm{\Phi}:\Omega\to\Omega^{\prime}$ 为 $k$ 阶微分同胚. 对分别定义在 $\Omega^{\prime}$ 和 $\Omega$ 上的函数 $u,v$, 用下式定义拉回算子 $\bm{\Phi}^*$ 和 $(\bm{\Phi}^{-1})^*$:
\begin{equation}\label{eq:sobolev-diffeomorphism-pullback-definitions}
\bm{\Phi}^*u=u\circ\bm{\Phi},\quad(\bm{\Phi}^{-1})^*v=v\circ\bm{\Phi}^{-1}
\end{equation}

\end{Definition}

\begin{Theorem}
\label{thm:sobolev-integer-order-pullback}
设 $\bm{\Phi}:\Omega\to\Omega^{\prime}$ 为 $k$ 阶微分同胚. 那么 $\bm{\Phi}^*$ 和 $(\bm{\Phi}^{-1})^*$ 分别为从 $H^k(\Omega^{\prime})$ 到 $H^k(\Omega)$ 以及从 $H^k(\Omega)$ 到 $H^k(\Omega^{\prime})$ 的有界线性映射.
\end{Theorem}

\begingroup
\renewcommand{\endProof}{\par}
\begin{Proof}
由于 $k$ 阶微分同胚的逆也是 $k$ 阶微分同胚, 只需对 $\bm{\Phi}^*$ 证明断言. 容易看出 $\bm{\Phi}^*$ 为从 $\widetilde L^2(\Omega^{\prime})$ 到 $\widetilde L^2(\Omega)$ 的有界线性映射, 因而由连续性也是从 $L^2(\Omega^{\prime})$ 到 $L^2(\Omega)$ 的有界线性映射. 注意 $(u,v)_{\Omega^{\prime}}=(\bm{\Phi}^*u,(\det\nabla\bm{\Phi})\bm{\Phi}^*v)_{\Omega}$.

还需考察导数. 为此, 需要证明 $\bm{\Phi}^*u$ 的导数可按通常方式利用链式法则和乘积法则计算. 若 $u\in C^{\infty}(\Omega)$, 这显然成立, 并且已经证明 $C^{\infty}$ 函数在 $H^k(\Omega)$ 中稠密. 对 $u\in H^k(\Omega)\cap C^{\infty}(\Omega)$, $\bm{\Phi}^*u$ 的任意第 $p$ 阶导数 $(p\leq k)$ 都是涉及 $u$ 的 $\leq p$ 阶导数的线性表达式, 其系数依赖于 $\bm{\Phi}$ 的导数; 特别地, 所有系数都是有界连续函数. 为将结果延拓到整个 $H^k(\Omega)$, 需要能够取极限. 以下引理表明可以这样做.
\end{Proof}
\endgroup

\begin{Lemma}
\label{lem:l-two-bounded-continuous-multiplier}
设 $u\in L^2(\Omega)\cap C(\Omega)$, $a\in C_b(\Omega)$. 那么 $au\in L^2(\Omega)$, 且 $\|au\|_2\leq\|a\|_{\infty}\|u\|_2$.
\end{Lemma}

\begingroup
\postdisplaypenalty=10000
\renewcommand{\endProof}{\par\nobreak\hfill$\blacksquare$\par\vspace{0.1cm}}
\begingroup
\renewenvironment{Proof}{\par}{\par\nobreak\hfill$\blacksquare$\par\vspace{0.1cm}}
\begin{Proof}
略去该引理的显然证明. 由于 $L^2(\Omega)\cap C(\Omega)$ 在 $L^2(\Omega)$ 中稠密, 该引理使我们可以对每个 $u\in L^2(\Omega)$ 定义 $au$. 至此完成 Theorem~\ref{thm:sobolev-integer-order-pullback} 的证明.
\end{Proof}
\endgroup
\endgroup
